% TEMPLATE-MILC-v3.tex - plantilla maestra para todas las guias MILC v3
% Ver scripts/generadores/build-guia-milc.py para la lista completa de
% claves esperadas. El script valida que no queden placeholders sin
% reemplazar antes de escribir el archivo final.

\documentclass[11pt,letterpaper]{article}
\usepackage[margin=1.75cm]{geometry}
\usepackage[table]{xcolor}
\usepackage{fontspec}
\usepackage[spanish]{babel}
\usepackage{tikz}
\usepackage[most]{tcolorbox}
\usepackage{tabularx}
\usepackage{array}
\usepackage{fancyhdr}
\usepackage{titlesec}
\usepackage{ragged2e}
\usepackage{enumitem}
\usepackage{microtype}

\setmainfont{Helvetica Neue}
\setsansfont{Helvetica Neue}
\setlength{\parindent}{0pt}
\setlength{\parskip}{6pt}
\hyphenpenalty=200
\exhyphenpenalty=200
\pretolerance=400
\tolerance=2000
\setlength{\emergencystretch}{1em}
\renewcommand{\arraystretch}{1.25}
\setlist{leftmargin=5mm,itemsep=1mm,topsep=1mm}
\newcolumntype{Y}{>{\RaggedRight\arraybackslash}X}

\definecolor{milcMagenta}{HTML}{E6007E}
\definecolor{milcVino}{HTML}{5A0038}
\definecolor{milcTurquesa}{HTML}{0093A5}
\definecolor{milcVerde}{HTML}{58A923}
\definecolor{milcMostaza}{HTML}{E5B400}
\definecolor{milcOcre}{HTML}{B86600}
\definecolor{milcNegro}{HTML}{241816}
\definecolor{milcGris}{HTML}{F7F7F4}
\definecolor{milcLinea}{HTML}{E8E2DD}
\definecolor{milcGradePrimary}{HTML}{84CC16}
\definecolor{milcGradeDark}{HTML}{4D7C0F}
\definecolor{milcGradeSoft}{HTML}{ECFCCB}
\definecolor{milcGradeAccent}{HTML}{CFE7FF}
\definecolor{milcGradeText}{HTML}{FFFFFF}
\color{milcNegro}

\pagestyle{fancy}
\fancyhf{}
\lhead{\small Tecnología e Informática · Bebras}
\rhead{\small Momento 3 · MILC v3}
\cfoot{\small\thepage}
\renewcommand{\headrulewidth}{0pt}

\newtcolorbox{titlebox}[1]{
  enhanced,colback=#1,colframe=#1,arc=7mm,boxrule=0pt,
  left=7mm,right=7mm,top=3mm,bottom=3mm,
  before skip=8pt,after skip=8pt,
  fontupper=\sffamily\bfseries\Large\color{white}
}

\newtcolorbox{softbox}[3]{
  enhanced,colback=#2,colframe=#1,borderline west={4pt}{0pt}{#1},
  arc=2mm,boxrule=.4pt,left=4mm,right=4mm,top=3mm,bottom=3mm,
  before skip=6pt,after skip=8pt,title={#3},coltitle=white,
  colbacktitle=#1,fonttitle=\sffamily\bfseries,fontupper=\RaggedRight,
  attach boxed title to top left={xshift=3mm,yshift=-2mm},
  boxed title style={arc=2mm, boxrule=0pt}
}

\newtcolorbox{drawbox}[2]{
  enhanced,colback=white,colframe=#1,arc=4mm,boxrule=1pt,
  left=5mm,right=5mm,top=4mm,bottom=4mm,before skip=8pt,after skip=8pt,
  title={#2},coltitle=white,colbacktitle=#1,fonttitle=\sffamily\bfseries,
  fontupper=\RaggedRight,
  attach boxed title to top left={xshift=4mm,yshift=-2mm},
  boxed title style={arc=3mm, boxrule=0pt}
}

\newtcolorbox{infoband}[1]{
  enhanced,colback=#1!8,colframe=#1!50,arc=2mm,boxrule=.5pt,
  left=4mm,right=4mm,top=2mm,bottom=2mm,before skip=4pt,after skip=4pt,
  fontupper=\small\RaggedRight
}

% Puente narrativo entre secciones (contrato editorial Conéctate).
% Da continuidad: "Acabas de X. Ahora vas a Y porque Z."
% Visualmente: banda izquierda en color del grado, texto en itálica pequeña.
\newtcolorbox{puentebox}{
  enhanced,colback=milcGradeSoft,colframe=milcGradeDark,
  borderline west={3pt}{0pt}{milcGradeDark},
  arc=0mm,boxrule=0pt,left=5mm,right=4mm,top=2.5mm,bottom=2.5mm,
  before skip=4pt,after skip=8pt,
  fontupper=\itshape\small\color{milcNegro}\RaggedRight
}
\newcommand{\puente}[1]{%
  \begin{puentebox}%
    \textcolor{milcGradeDark}{$\rightarrow$}\hspace{2mm}#1%
  \end{puentebox}%
}

\newcommand{\linead}[1]{\noindent\color{milcLinea}\rule{#1}{0.5pt}\color{milcNegro}}
\newcommand{\respuesta}[1]{\par\vspace{2mm}\foreach \n in {1,...,#1}{\noindent\color{milcLinea}\rule{\linewidth}{0.5pt}\par\vspace{5mm}}\color{milcNegro}}
\newcommand{\checkbox}{\textcolor{milcGradeDark}{\fbox{\phantom{x}}}\hspace{5pt}}

\newcommand{\phasecard}[4]{
\begin{tcolorbox}[enhanced,colback=#3,colframe=#2,arc=3mm,boxrule=.5pt,height=3.05cm,valign=center,halign=center,left=1.5mm,right=1.5mm,top=2mm,bottom=2mm]
{\sffamily\bfseries\fontsize{22}{24}\selectfont\textcolor{#2}{#1}}\par
{\sffamily\bfseries\small #4}
\end{tcolorbox}
}

\begin{document}
\thispagestyle{empty}
\begin{tikzpicture}[remember picture,overlay]
  \fill[white] (current page.south west) rectangle (current page.north east);
  \path[fill=milcGradePrimary,rounded corners=16mm]
    ([xshift=.55cm,yshift=-.55cm]current page.north west) rectangle
    ([xshift=-.55cm,yshift=3.65cm]current page.south east);

  \node[anchor=north west,text=milcGradeText] at ([xshift=2.0cm,yshift=-1.85cm]current page.north west)
    {\sffamily\bfseries\fontsize{14}{16}\selectfont MOMENTO};
  \node[anchor=north west,text=milcGradeText,opacity=.82] at ([xshift=2.0cm,yshift=-2.75cm]current page.north west)
    {\sffamily\bfseries\fontsize{9}{11}\selectfont BEBRAS · PENSAR COMO UNA MÁQUINA};

  \begin{scope}[shift={([xshift=-3.05cm,yshift=-2.55cm]current page.north east)}]
    \fill[white,opacity=.80] (-.62,-.52) rectangle (.62,.52);
    \draw[milcGradeText,line width=1pt] (-.62,-.52) rectangle (.62,.52);
    \fill[milcGradeDark] (-.42,-.38) rectangle (-.22,.06);
    \fill[milcGradeAccent] (-.10,-.38) rectangle (.10,.24);
    \fill[milcGradePrimary!60!black] (.22,-.38) rectangle (.42,.38);
  \end{scope}

  \node[anchor=west,text=milcGradeText] at ([xshift=1.9cm,yshift=-12.85cm]current page.north west)
    {\sffamily\bfseries\fontsize{124}{124}\selectfont 3};
  \draw[milcGradeText,line width=1.7pt]
    ([xshift=4.52cm,yshift=-11.68cm]current page.north west) circle (.13cm);

  \node[anchor=west,text=milcGradeText,text width=8.7cm,align=left] at ([xshift=10.35cm,yshift=-11.6cm]current page.north west)
    {
     {\sffamily\bfseries\fontsize{19}{22}\selectfont Lo esencial y nada más --- la abstracción\\como tercera pieza del pensamiento computacional}\\[4mm]
     {\sffamily\bfseries\fontsize{23}{26}\selectfont Bebras · Momento 3}\\[2mm]
     {\sffamily\fontsize{13}{16}\selectfont Curso «Pensar como una máquina» · Pensamiento computacional}\\[5mm]
     {\sffamily\bfseries\fontsize{11}{13}\selectfont Institución Educativa Sor María Juliana}\\[1mm]
     {\sffamily\fontsize{10}{12}\selectfont Autor: Dr. Álvaro Cárdenas Orozco}
    };

  \path[fill=milcGradeSoft,rounded corners=7mm]
    ([xshift=1.35cm,yshift=.85cm]current page.south west) rectangle
    ([xshift=-1.35cm,yshift=4.25cm]current page.south east);
  \draw[milcGradeDark,line width=.65pt,rounded corners=7mm]
    ([xshift=1.35cm,yshift=.85cm]current page.south west) rectangle
    ([xshift=-1.35cm,yshift=4.25cm]current page.south east);
  \node[anchor=center,text=milcNegro,text width=17.0cm,align=center] at ([yshift=2.55cm]current page.south)
    {\sffamily\fontsize{10}{12}\selectfont Metodología MILC · Apertura ancestral · Pensamiento computacional · Triángulo Dussel-Estoicismo-Floridi\\
    \textcolor{milcGradeDark}{\bfseries Producto:} Tu modelo escrito de algo complejo (conservo · ignoro) más un mapa o ícono a mano alzada que abstraiga un lugar de tu barrio, con leyenda.};
\end{tikzpicture}
\mbox{}
\newpage

\section*{Apertura: Lo esencial y nada más: la abstracción como tercera pieza del pensamiento computacional}

\begin{softbox}{milcGradeDark}{milcGradeSoft}{Saber ancestral}
En las veredas del Valle del Cauca, cuando un campesino te explica cómo
llegar a su finca, no describe cada árbol del cafetal ni el color de cada
casa del camino. Te da solo lo esencial: \textbf{«pasas la tienda de doña
Rosa, coges la trocha de la derecha y en el guayacán grande volteas»}. Tres
señas y ya sabes llegar. No menciona el perro que ladra en la curva, ni
cuántos postes de luz hay, ni el nombre del vecino, porque nada de eso te
ayuda a no perderte. Ese campesino hace, sin nombrarlo, un trabajo
finísimo: \textbf{separa las pocas señas que de verdad importan} del montón
de detalles que solo serían ruido. Limpiar el camino hasta dejar lo justo
--- ni una seña de más, ni una de menos --- es abstracción pura, heredada
en la tradición oral del campo mucho antes de que existieran los mapas
digitales.
\end{softbox}

\begin{softbox}{milcTurquesa}{milcGris}{Saber contemporáneo}
\textbf{Abstraer} es quedarte con lo esencial e ignorar lo que no importa
para la tarea que tienes entre manos. Es construir un \textbf{modelo}: una
versión simplificada de la realidad que conserva solo los datos que
cambian el resultado. Un mapa de Cartago no dibuja cada ladrillo ni cada
ventana; dibuja las calles y unos puntos clave. Una app de clima no
almacena el nombre de cada nube; guarda temperatura, humedad y presión.
Un ícono de batería muestra el nivel, no la química interna. En todos los
casos la idea es la misma: \textbf{menos ruido es más claridad}. En Bebras,
muchas tareas traen datos de sobra a propósito; abstraer es saber
\emph{cuáles} importan para la pregunta y cuáles son puro ruido para
distraerte.
\end{softbox}

\begin{softbox}{milcMostaza}{milcGris}{Pregunta-puente}
\textit{El campesino abstrae el camino a su vereda: ignora el ruido y conserva
solo lo que sirve para llegar. ¿Y si te dijera que un computador hace
exactamente eso con cada modelo que construye --- y que en muchas tareas
Bebras lo más difícil no es calcular, sino decidir \emph{cuáles} datos
importan?}
\end{softbox}

\begin{softbox}{milcVerde}{milcGris}{Saber y saber hacer}
Al terminar podrás: (1) \textbf{identificar} qué es una abstracción y
reconocer modelos cotidianos como mapas, íconos y resúmenes; (2)
\textbf{explicar} con tus palabras qué conservar y qué ignorar al
modelar algo complejo de tu entorno; (3) \textbf{aplicar} la abstracción
para ignorar datos distractores en un reto estilo Bebras y llegar a la
respuesta correcta.
\end{softbox}

\small
\begin{tabularx}{\textwidth}{>{\bfseries}p{3.0cm}Y}
\rowcolor{milcGradePrimary!12}Autor & Dr. Álvaro Cárdenas Orozco\\
DBA & Desarrolla el pensamiento computacional --- descomposición, patrones, abstracción y algoritmos --- para resolver problemas (transversal 6°--11°, alineado a los lineamientos MEN de Tecnología e Informática)\\
Referentes & Valentina Dagienė (Bebras) · Pensamiento computacional (Wing) · Dussel · Estoicismo · Floridi · Saberes del Valle y el Pacífico\\
Producto final & Tu modelo escrito de algo complejo (conservo · ignoro) más un mapa o ícono a mano alzada que abstraiga un lugar de tu barrio, con leyenda.\\
\end{tabularx}
\normalsize

\puente{Ya sabes que el campesino reduce el camino a tres señas. Ahora mira el
mapa de la sesión: vas a recorrer los mismos movimientos de la
abstracción hasta modelar algo tuyo y resolver un reto con datos
distractores.}

\begin{titlebox}{milcGradeDark}
Ruta MILC: así vamos a aprender
\end{titlebox}
\begin{tabularx}{\textwidth}{>{\centering\arraybackslash}X>{\centering\arraybackslash}X>{\centering\arraybackslash}X>{\centering\arraybackslash}X}
\phasecard{1}{milcVerde}{milcGris}{Escucha\\[-1mm]\small Observo, escucho y pregunto.} &
\phasecard{2}{milcTurquesa}{milcGris}{Sistematización\\[-1mm]\small Pensamiento computacional.} &
\phasecard{3}{milcMagenta}{milcGris}{Praxis\\[-1mm]\small Construyo y aplico.} &
\phasecard{4}{milcVino}{milcGris}{Evaluación\\[-1mm]\small Reflexión liberadora.}
\end{tabularx}


\puente{Antes de nombrar la pieza, conviene verla en tu propia experiencia.
Recuerda una vez en que alguien te dio indicaciones para llegar a algún
lugar: la abstracción ya estaba ahí, aunque nadie le pusiera ese nombre.}

\begin{titlebox}{milcVerde}
Fase 1 · Escucha
\end{titlebox}

\begin{softbox}{milcVerde}{milcGris}{Escena para escuchar}
\textbf{Recuerda una vez en que alguien te dio indicaciones para llegar a
un lugar} --- un familiar, un vecino, un amigo. Puede ser cómo llegar a
una finca, a una casa nueva o a un lugar del barrio. Fíjate en qué
\emph{sí} te dijeron y qué \emph{no} te dijeron. ¿Mencionaron el color
de cada casa? ¿El nombre de cada árbol? ¿O solo las señas que de verdad
te servían para no perderte? Cada dato que dejaron por fuera fue una
decisión: ese detalle no cambiaba el resultado. Sin saberlo, la persona
que te dio las indicaciones ya estaba abstrayendo.
\end{softbox}

\checkbox Recordaste al menos tres señas o datos que sí te dieron para llegar
(lo que conservaron).\\
\checkbox Nombraste al menos dos detalles que nadie mencionó, aunque existían
en el camino (lo que ignoraron).\\
\checkbox Explicaste, en una frase, por qué esos detalles ignorados no cambiaban
el resultado de llegar.

\begin{infoband}{milcVerde}
\textbf{En tu cuaderno (Actividad 1 · IDENTIFICA).} Título: «Actividad 1
--- Las señas que importan». \textbf{Formato:} dos listas cortas:
«conservaron» e «ignoraron». \textbf{Extensión:} 3 ítems en la primera
lista y 2 en la segunda, más una frase final que explique por qué lo
ignorado no importaba. \textbf{Sabes que terminaste cuando} tu lista
muestra claramente qué se quedó y qué se descartó, y por qué.
\end{infoband}

\puente{Acabas de ver la abstracción en lo cotidiano. Ahora vamos a nombrarla y
entender qué conserva y qué descarta un buen modelo, porque con eso vas a
resolver todo lo que sigue.}

\begin{titlebox}{milcTurquesa}
Fase 2 · Sistematización con pensamiento computacional
\end{titlebox}

\begin{softbox}{milcTurquesa}{milcGris}{Cuatro pilares aplicados al tema}
Abstraer es la tercera pieza del pensamiento computacional. Su trabajo es
\textbf{construir un modelo}: una versión simplificada de algo real que
conserva solo lo que importa para la tarea. No es resumir mal --- es
\emph{elegir} a propósito qué se queda y qué se va. Aquí están las cuatro
ideas que hacen funcionar esta pieza.
\end{softbox}

\small
\begin{tabularx}{\textwidth}{>{\bfseries\RaggedRight\arraybackslash}p{3.6cm}Y}
\rowcolor{milcTurquesa!90}\color{white}Pilar & \color{white}\bfseries Aplicación al tema\\
1. Descomposición & \textbf{Lo esencial es lo que cambia el resultado.} Un dato es esencial si,
al quitarlo, la respuesta cambia. Si al quitarlo la respuesta es la misma,
es ruido. Las señas de la vereda son esenciales; el color de la casa de al
lado, no.\\
2. Reconocimiento de patrones & \textbf{Un modelo nunca es la realidad: es lo suficiente.} Un mapa de
Cartago que dibuja cada ladrillo no sirve para guiarse: hay demasiado ruido.
El mejor modelo no es el más completo, sino el que tiene \emph{los datos
justos} para la tarea.\\
3. Abstracción & \textbf{Abstraer es una decisión, no un accidente.} El campesino elige
cuáles señas dar. El diseñador elige qué poner en el ícono. Tú decides
qué datos del enunciado Bebras importan. Cada vez que abstraes estás
tomando una decisión sobre qué parte de la realidad vale la pena ver.\\
4. Algoritmo conceptual & \textbf{Menos ruido es más claridad.} Cuando quitas los detalles que
estorban, tu cabeza puede concentrarse en lo que de verdad mueve la aguja.
Por eso abstraer no es pereza: es una habilidad que hace el pensamiento
más rápido y más preciso.\\
\end{tabularx}
\normalsize

\begin{drawbox}{milcTurquesa}{Cómo abstraer un problema Bebras}
\textbf{1. Lee el enunciado} sin afán y subraya los datos que recibes. \\[1mm]
\textbf{2. Marca la pregunta} exacta que te hacen. \\[1mm]
\textbf{3. Clasifica} cada dato: ¿cambia la respuesta si lo quito? Si sí,
es esencial. Si no, es ruido. \\[1mm]
\textbf{4. Tacha el ruido} y trabaja solo con los datos esenciales. \\[1mm]
\textbf{5. Responde} con lo que quedó. Si te trabas, revisa si hay más
ruido que no descartaste.
\end{drawbox}

\begin{softbox}{milcMostaza}{milcGris}{Errores comunes y su corrección}
\textbf{Confundir abstraer con resumir mal:} abstraer es \emph{elegir} qué
dato importa, no quitar al azar. \textbf{Usar todos los datos del enunciado:}
muchas tareas Bebras traen ruido a propósito; usar todo casi siempre
complica la respuesta. \textbf{Confundir esencial con «el número más
grande»:} el dato esencial es el que cambia la respuesta, no el que parece
más importante a primera vista.
\end{softbox}

\begin{infoband}{milcTurquesa}
\textbf{En tu cuaderno (Actividad 2 · EXPLICA).} Título: «Actividad 2 ---
Los cuatro movimientos de la abstracción». \textbf{Formato:} tabla de 2
columnas (movimiento · ejemplo de mi día). \textbf{Extensión:} los cuatro
movimientos, cada uno con un ejemplo distinto a los del texto.
\textbf{Sabes que terminaste cuando} un compañero entendería tus ejemplos
sin que se los expliques.
\end{infoband}

\puente{Ya distingues lo esencial del ruido. Es hora de usarlo: vas a modelar
algo complejo de tu entorno y luego ignorar datos distractores en un
reto estilo Bebras.}

\begin{titlebox}{milcMagenta}
Fase 3 · Praxis con pensamiento computacional
\end{titlebox}

\begin{softbox}{milcMagenta}{milcGris}{Construyendo el producto con los cuatro pilares}
La abstracción no es adorno: se usa para resolver. Vas a enfrentarte a un
\textbf{reto estilo Bebras con datos de sobra}, igual que en el concurso
real. La clave no es calcular rápido --- es decidir primero qué datos
importan y cuáles son puro ruido.
\end{softbox}

\small
\begin{tabularx}{\textwidth}{>{\bfseries\RaggedRight\arraybackslash}p{3.6cm}Y}
\rowcolor{milcMagenta!90}\color{white}Pilar & \color{white}\bfseries Aplicación al producto\\
Descomposición & \textbf{Abstracción al foco:} antes de calcular nada, identifica la
pregunta exacta. Solo cuando sabes qué te piden puedes decidir qué dato
importa.\\
Patrones & \textbf{Clasifica los datos:} recorre cada dato del enunciado y pregúntate:
«¿si lo quito, cambia la respuesta?». Separa los esenciales de los que
son ruido.\\
Abstracción & \textbf{Descarta con decisión:} tacha el ruido en tu mente o en el papel.
No lo uses. Tener datos que no sirven y no usarlos es tan importante como
usar los que sí sirven.\\
Algoritmo de armado & \textbf{Responde con lo esencial:} trabaja solo con los datos que
sobrevivieron al descarte. Si la respuesta parece rara, revisa si
descartaste algo que sí importaba o si usaste algo que era ruido.\\
\end{tabularx}
\normalsize

\begin{drawbox}{milcMagenta}{El reto: la tienda de la esquina}
\checkbox \textbf{Reto (Nivel C).} En una tienda de Cartago, una gaseosa
cuesta \textbf{\$2.500}. El local tiene \textbf{7 sabores} y abre a las
\textbf{8~a.m.} Compras \textbf{3 gaseosas}. ¿Cuánto pagas en total? \\[2mm]
\checkbox Marca: \quad \fbox{\phantom{x}}~\$5.000 \quad
\fbox{\phantom{x}}~\$7.500 \quad \fbox{\phantom{x}}~\$10.000 \quad
\fbox{\phantom{x}}~\$17.500 \\[2mm]
\checkbox Escribe qué datos eran \textbf{ruido} y por qué los ignoraste.
\end{drawbox}

\begin{softbox}{milcMagenta}{milcGris}{Plantilla de trabajo}
\textbf{Resuélvelo paso a paso (abstracción + algoritmo).} La pregunta es
cuánto pagas por 3 gaseosas. Para responderla solo necesitas: el precio
(\$2.500) y la cantidad (3). Los 7 sabores y la hora de apertura (8~a.m.)
no entran en el cálculo --- son ruido a propósito, para ver si los
descartas. Operación: $3 \times \$2.500 = \$7.500$. El ruido eran el
número de sabores y la hora de apertura; ninguno cambia cuánto pagas.
\end{softbox}

\begin{infoband}{milcMagenta}
\textbf{En tu cuaderno (Actividad 3 · APLICA).} Título: «Actividad 3 ---
Mi primer reto Bebras de abstracción». \textbf{Formato:} respuesta marcada
+ lista de datos de ruido + 2 renglones de explicación. \textbf{Extensión:}
la respuesta, los datos que descartaste nombrados, y la razón en una frase.
\textbf{Sabes que terminaste cuando} tu explicación convencería a alguien
que aún no sabe la respuesta, diciéndole qué ignoró y por qué.
\end{infoband}

\puente{Resolviste el reto. Ahora deja tu evidencia: el modelo escrito y el mapa
o ícono que muestran que ya sabes quedarte con lo esencial.}

\begin{titlebox}{milcMostaza}
Producto de la sesión
\end{titlebox}
\textbf{Modelo escrito (conservo · ignoro) + mapa o ícono de tu barrio con leyenda}.
\begin{softbox}{milcMostaza}{milcGris}{Criterios de calidad}
(1) Tu modelo escrito nombra al menos 3 datos esenciales y 2 datos de ruido
del objeto o lugar que elegiste. (2) Explicaste en una frase por qué cada dato de ruido no cambia el resultado. (3) Tu mapa o ícono deja por fuera los detalles que no ayudan y conserva
solo lo esencial para entender el lugar. (4) La leyenda del mapa o ícono explica qué representa cada seña o símbolo. (5) En el reto de la tienda marcaste \$7.500 y nombraste los dos datos de
ruido (los 7 sabores y la hora de apertura).
\end{softbox}

\puente{Terminaste el producto. Detente a mirar qué cambió en ti --- no la nota,
sino tu manera de ver la información que te rodea.}

\begin{titlebox}{milcVino}
Fase 4 · Evaluación liberadora — cinco dimensiones
\end{titlebox}

\begin{softbox}{milcVino}{milcGris}{Sentido de la evaluación}
Evaluar no es solo revisar una nota. Es reconocer qué aprendí, qué cambió en mí, qué emociones manejé, cómo me relaciono con la ciudad y la comunidad, y qué le dejaré a quien viene detrás.
\end{softbox}

\textbf{Mi reflexión liberadora en cinco dimensiones.} Completa cada frase iniciada:

\small
\begin{tabularx}{\textwidth}{>{\bfseries\RaggedRight\arraybackslash}p{3.4cm}Y>{\RaggedRight\arraybackslash}p{4.6cm}}
\rowcolor{milcVino}\color{white}Dimensión & \color{white}\bfseries Mi respuesta (frame starter) & \color{white}\bfseries Algo que descubrí\\
1. Personal & Lo que más cambió en mí fue \linead{6.5cm} & \linead{4.2cm}\\
2. Emocional & La emoción más fuerte fue \linead{2.5cm} y la regulé \linead{3cm} & \linead{4.2cm}\\
3. Ciudadana & En democracia esto importa porque \linead{6cm} & \linead{4.2cm}\\
4. Local (Cartago) & En mi barrio sería útil para \linead{6.5cm} & \linead{4.2cm}\\
5. Intergeneracional & A mi abuela/abuelo le diría \linead{6.5cm} & \linead{4.2cm}\\
\end{tabularx}
\normalsize

\begin{infoband}{milcVino}
\textbf{Para la próxima sustentación.} Vuelve a esta tabla en seis meses. Verás que las respuestas cambian — no porque lo aprendido cambie, sino porque tú cambias. La evaluación liberadora es una conversación contigo a través del tiempo.
\end{infoband}


\puente{Cerramos pensando en grande: tres voces para entender qué significa, para
ti y para tu comunidad, aprender a limpiar el ruido del mundo.}

\begin{titlebox}{milcGradeDark}
Triángulo de pensamiento · cierre filosófico
\end{titlebox}

\begin{softbox}{milcVino}{milcGris}{Dussel · lente del nosotros}
\textit{``El saber del pueblo también es ciencia.''}

Dar buenas señas para llegar a una vereda ya era abstraer: el campesino
limpiaba el camino hasta dejar lo esencial, mucho antes de que existieran
los mapas digitales. Ese saber no llegó de afuera; tu comunidad ya lo
practicaba, aunque nadie lo llamara pensamiento computacional.

\textbf{¿Qué saberes de mi familia o mi vereda ya abstraían lo esencial, aunque nadie los llamara abstracción?}
\end{softbox}
\linead{\linewidth}\\[1mm]
\linead{\linewidth}

\begin{softbox}{milcMostaza}{milcGris}{Estoicismo · Marco Aurelio · cuidado interior}
\textit{``En cada acto pregúntate: ¿es esto necesario? Suprime lo superfluo, y ganarás tiempo y serenidad.''}

Abstraer un problema es un acto estoico: quitas lo que estorba y te
quedas con lo que importa. Lo mismo sirve fuera del papel: cuánto ruido
--- notificaciones, preocupaciones ajenas, distracciones --- lleva tu
atención lejos de lo esencial.

\textbf{Igual que abstraigo un problema quitándole el ruido, ¿qué «ruido» de mi vida podría empezar a ignorar para enfocarme en lo que de verdad importa?}
\end{softbox}
\linead{\linewidth}\\[1mm]
\linead{\linewidth}

\begin{softbox}{milcTurquesa}{milcGris}{Floridi · lente de la infoesfera}
\textit{``Todo modelo es un nivel de abstracción: muestra unas cosas y, al hacerlo, deja otras fuera.''}

Un mapa, un ícono, un perfil de red social: todos son abstracciones que
deciden qué parte de la realidad se ve y qué queda invisible. Esa
decisión no es neutral: quien construye el modelo elige qué importa.

\textbf{Si un mapa o un perfil son abstracciones de algo más grande, ¿qué se pierde cada vez que abstraemos — y quién decide qué parte de la realidad queda por fuera?}
\end{softbox}
\linead{\linewidth}\\[1mm]
\linead{\linewidth}

\begin{titlebox}{milcVino}
Compromiso final
\end{titlebox}

\small
\begin{tabularx}{\textwidth}{>{\RaggedRight\arraybackslash}p{3.0cm}YYY}
\rowcolor{milcVino}\color{white}\bfseries Criterio & \color{white}\bfseries Lo logré & \color{white}\bfseries En proceso & \color{white}\bfseries Necesito apoyo\\
Comprensión del tema & Explico ideas centrales con mis palabras. & Reconozco ideas pero falta precisión. & Necesito repasar conceptos base.\\
Pensamiento computacional & Apliqué los 4 pilares al diseñar mi producto. & Apliqué algunos pilares. & Necesito releer la fase 2.\\
Aplicación y producto & Mi producto cumple los criterios y se puede revisar. & Producto funcional con ajustes pendientes. & Aún en construcción.\\
Reflexión 5 dimensiones & Respondí cada dimensión con honestidad. & Respondí algunas con detalle. & Mis respuestas son muy generales.\\
Triángulo de pensamiento & Conecté las 3 voces con mi trabajo real. & Conecté al menos una. & No relaciono las voces con mi proceso.\\
\end{tabularx}
\normalsize

\textbf{Hoy entendí que:}\\[1mm]
\linead{\linewidth}\\[1mm]
\linead{\linewidth}

\textbf{Mi compromiso para la próxima sesión es:}\\[1mm]
\linead{\linewidth}\\[1mm]
\linead{\linewidth}

\begin{softbox}{milcMostaza}{milcGris}{Cita de despedida · Marco Aurelio}
\textit{``El obstáculo en el camino se vuelve el camino. Lo que estorba al camino se vuelve camino.''}\\[1mm]
Cada error de hoy, bien observado, se convierte en pieza del camino que pisarás mañana.
\end{softbox}

\small
\begin{tabularx}{\textwidth}{>{\bfseries}p{3.6cm}Y}
\rowcolor{milcGradeDark!12}Nombre completo: & \linead{10cm}\\
Fecha: & \linead{4cm} \quad Curso/grupo: \linead{4cm}\\
Mi firma: & \linead{9cm}\\
\end{tabularx}
\normalsize

\begin{infoband}{milcGradeDark}
\textbf{Para la próxima sesión.} Elige una app, una señal de tránsito o un
ícono que uses con frecuencia --- puede ser el ícono del wifi, la señal de
«pare» o el logo de tu tienda favorita --- y escribe en tu cuaderno: ¿qué
conserva esa abstracción? ¿Qué dejó por fuera? ¿Cambiaría algo si hubiera
guardado un detalle más? Trae tu ejemplo para compartirlo con el grupo.
\end{infoband}

\end{document}
